\section{Conclusion}

This work delivered an empirical benchmark of smart contract security tooling over 30 Solidity contracts organized as Vulnerable/Fixed MVE pairs across the OWASP Smart Contract Top 10. The Vuln/Fixed design exposed an asymmetry that aggregated benchmarks tend to hide. Static and symbolic tools achieved lower recall than fuzzers (Slither 7/15, Mythril 6/15, against 10/15 for both Echidna and Medusa) and at substantially higher false-positive cost (15/15 and 10/15 respectively, against 0/15 for fuzzing). Fuzzers dominated this corpus on both axes; the question for an audit pipeline is therefore not whether to use them, but how to compose them with paradigms that capture what they cannot.

The data support a minimal three-stage audit pipeline: Slither $\rightarrow$ Mythril $\rightarrow$ Echidna. Slither operates as a fast first-pass filter on suspicious patterns; Mythril provides path-level confirmation on flagged functions through symbolic execution over EVM bytecode; Echidna validates that residual invariants resist tens of thousands of randomized transactions. Halmos extends this stack with formal proof on categories where \texttt{check\_*} properties can be authored — demonstrated here on SC01 and SC09, where the property authored against the Vulnerable contract produced a symbolic counterexample and the equivalent property against the Fixed contract was proven over the entire symbolic input space. Unlike a passing fuzzer or a silent static analyzer, this is conclusive evidence that no input exists for which the property fails.

Aderyn contributed 3/15 detections at the same 15/15 false-positive rate as Slither, with positives that largely overlap Slither's coverage — limited additive value beyond the stronger semantic analyzer. Semgrep and Solhint produced no exploit-level findings and serve only as standard-tooling baselines: their 0/15 detection rate measures what protection the default Solidity CI pipeline offers without dedicated security tooling. LLM review is useful for triage and qualitative explanation, but the gap between the best and worst calibration in this benchmark ($\Delta$ +3.3 for LLaMA-3.3-70B against +1.4 for Gemini 2.5 Flash) shows that the paradigm is sensitive enough to model and prompt that it cannot serve as a verdict layer in an automated pipeline. The structural absence of executable grounding remains the deeper reason.

Two limitations bound the empirical reach of this benchmark. First, the MVE design produces compact, single-vector contracts: each Fixed/Vulnerable pair isolates one OWASP category, which simplifies attribution but excludes interaction effects between vectors that real-world contracts exhibit. Second, the fuzzing harness for SC05 was intentionally minimal — generic invariant functions without ECDSA generation or attacker contracts — while the SC08 harnesses do include dedicated attacker contracts; the false negatives observed on SC05 therefore reflect harness expressiveness, not a structural limit of property-based fuzzing itself. Three extensions would broaden the result. Authoring \texttt{check\_*} properties for the remaining eight OWASP categories would convert Halmos from per-category formal proof to a coverage-wide guarantee. Replacing or augmenting the synthetic MVE corpus with real-world contracts would test whether tool behavior generalizes to production code. Calibrating the LLM prompt against a dedicated mitigation-recognition test set could close the precision gap observed on Gemini, likely a prompt-engineering problem rather than a model-capacity limitation, since LLaMA and Qwen achieved strong discrimination with the same prompt.